\documentclass[11pt]{article}
\usepackage[utf8]{inputenc}
\usepackage[T1]{fontenc}
\usepackage[spanish]{babel}
\usepackage{amsmath, amssymb, amsthm}
\usepackage{braket}
\usepackage[margin=2.5cm]{geometry}
\usepackage{hyperref}

\allowdisplaybreaks

\newcommand{\dd}{\mathrm{d}}
\newcommand{\ii}{\mathrm{i}}
\newcommand{\ee}{\mathrm{e}}
\newcommand{\vb}[1]{\mathbf{#1}}
\newcommand{\uv}[1]{\hat{\vb e}_{#1}}

\title{Cuantizaci\'on del campo electromagn\'etico}
\author{Ricardo Del Rosal Gardu\~no \\ \small Doctorado en F\'isica, BUAP}
\date{Notas a partir de Schleich, \emph{Quantum Optics in Phase Space}}

\begin{document}
\maketitle

\section{Leyes de Maxwell}

El punto de partida es el campo electromagn\'etico cl\'asico. La ley de Gauss
para el campo el\'ectrico afirma que el flujo de $\vb E$ a trav\'es de una
superficie cerrada es proporcional a la carga el\'ectrica encerrada por esa
superficie. En forma integral,
\begin{equation}
  \oint_S \vb E\cdot \dd\vb A
  =
  \frac{Q_{\mathrm{enc}}}{\varepsilon_0}.
\end{equation}
La forma diferencial correspondiente es
\begin{equation}
  \nabla\cdot \vb E(\vb r,t)
  =
  \frac{\rho(\vb r,t)}{\varepsilon_0}.
\end{equation}
Esta ecuaci\'on local dice que la densidad de carga es fuente del campo
el\'ectrico: donde hay carga positiva la divergencia de $\vb E$ es positiva,
y donde hay carga negativa es negativa.

Recordando la ley de Coulomb, la fuerza que ejerce una carga $q_1$ localizada
en $\vb r_1$ sobre una carga $q_2$ localizada en $\vb r_2$ es
\begin{equation}
  \vb F_{12}
  =
  \frac{1}{4\pi\varepsilon_0}
  q_1q_2
  \frac{\vb r_2-\vb r_1}{|\vb r_2-\vb r_1|^3}.
\end{equation}
Dividiendo por la carga de prueba se obtiene el campo el\'ectrico producido
por $q_1$:
\begin{equation}
  \vb E(\vb r)
  =
  \frac{\vb F}{q_0}
  =
  \frac{1}{4\pi\varepsilon_0}
  q_1
  \frac{\vb r-\vb r_1}{|\vb r-\vb r_1|^3}.
\end{equation}
Por superposici\'on lineal, para $N$ cargas puntuales,
\begin{equation}
  \vb E_{\mathrm{tot}}(\vb r)
  =
  \sum_{i=1}^N
  \frac{1}{4\pi\varepsilon_0}
  q_i
  \frac{\vb r-\vb r_i}{|\vb r-\vb r_i|^3}.
\end{equation}
Para una distribuci\'on continua de carga, esta suma se reemplaza por una
integral:
\begin{equation}
  \vb E(\vb r)
  =
  \frac{1}{4\pi\varepsilon_0}
  \int_V
  \rho(\vb r')
  \frac{\vb r-\vb r'}{|\vb r-\vb r'|^3}
  \dd^3r'.
\end{equation}

La forma integral de la ley de Gauss puede entenderse directamente para una
carga puntual. Si la carga $q$ est\'a dentro del volumen encerrado por $S$,
entonces el flujo es
\begin{equation}
  \Phi_E
  =
  \oint_S \vb E\cdot \dd\vb A
  =
  \frac{q}{4\pi\varepsilon_0}
  \oint_S
  \frac{\hat{\vb r}\cdot \dd\vb A}{r^2}.
\end{equation}
El cociente $(\hat{\vb r}\cdot \dd\vb A)/r^2$ es el elemento de \'angulo
s\'olido $\dd\Omega$ subtendido por el elemento de \'area. Por tanto,
\begin{equation}
  \Phi_E
  =
  \frac{q}{4\pi\varepsilon_0}
  \oint_S \dd\Omega.
\end{equation}
Si la carga est\'a dentro de $S$, la superficie cerrada cubre todo el \'angulo
s\'olido y $\oint_S\dd\Omega=4\pi$. Si la carga est\'a fuera, las
contribuciones entrantes y salientes se cancelan y el resultado es cero.
Por eso,
\begin{equation}
  \oint_S \vb E\cdot \dd\vb A
  =
  \frac{q}{\varepsilon_0}
  \quad\text{si }q\text{ est\'a dentro de }V.
\end{equation}
Con varias cargas se suma s\'olo sobre las cargas encerradas:
\begin{equation}
  \oint_S \vb E\cdot \dd\vb A
  =
  \frac{1}{\varepsilon_0}
  \sum_{i\in V}q_i.
\end{equation}
Para una distribuci\'on continua,
\begin{equation}
  \oint_S \vb E\cdot \dd\vb A
  =
  \frac{1}{\varepsilon_0}
  \int_V \rho(\vb r)\,\dd^3r.
\end{equation}
Aplicando el teorema de Gauss,
\begin{equation}
  \oint_S \vb E\cdot \dd\vb A
  =
  \int_V \nabla\cdot\vb E\,\dd^3r,
\end{equation}
se obtiene
\begin{equation}
  \int_V \nabla\cdot\vb E\,\dd^3r
  =
  \int_V \frac{\rho(\vb r)}{\varepsilon_0}\,\dd^3r.
\end{equation}
Como el volumen es arbitrario, los integrandos deben coincidir:
\begin{equation}
  \boxed{\nabla\cdot\vb E=\frac{\rho}{\varepsilon_0}.}
\end{equation}
Para una carga puntual en $\vb r_0$, la densidad se escribe como
\begin{equation}
  \rho(\vb r)=q\,\delta^{(3)}(\vb r-\vb r_0).
\end{equation}

La ley de Gauss para el campo magn\'etico establece que el flujo magn\'etico a
trav\'es de cualquier superficie cerrada es cero:
\begin{equation}
  \oint_S \vb B\cdot \dd\vb A=0,
  \qquad
  \boxed{\nabla\cdot\vb B=0.}
\end{equation}
Esta ecuaci\'on expresa la ausencia de monopolos magn\'eticos en la teor\'ia
electromagn\'etica usual. Las l\'ineas de campo magn\'etico no empiezan ni
terminan en cargas magn\'eticas; forman lazos cerrados o se extienden al
infinito.

La ley de Faraday-Lenz dice que un flujo magn\'etico variable induce un campo
el\'ectrico rotacional. En forma integral,
\begin{equation}
  \oint_C \vb E\cdot \dd\vb \ell
  =
  -\frac{\dd}{\dd t}
  \int_S \vb B\cdot \dd\vb A.
\end{equation}
Usando el teorema de Stokes se obtiene
\begin{equation}
  \boxed{\nabla\times\vb E=-\frac{\partial\vb B}{\partial t}.}
\end{equation}
El signo menos es la ley de Lenz: el campo inducido se opone al cambio de
flujo que lo produce.

Finalmente, la ley de Ampere-Maxwell establece que los campos magn\'eticos
pueden ser generados por corrientes el\'ectricas y por campos el\'ectricos que
var\'ian en el tiempo:
\begin{equation}
  \oint_C \vb B\cdot \dd\vb \ell
  =
  \mu_0 I_{\mathrm{enc}}
  +
  \mu_0\varepsilon_0
  \frac{\dd}{\dd t}
  \int_S \vb E\cdot \dd\vb A.
\end{equation}
Su forma diferencial es
\begin{equation}
  \boxed{
  \nabla\times\vb B
  =
  \mu_0\vb J
  +
  \mu_0\varepsilon_0
  \frac{\partial\vb E}{\partial t}.
  }
\end{equation}
En ausencia de cargas y corrientes, las ecuaciones de Maxwell se reducen a
las ecuaciones del campo de radiaci\'on libre.

\section{Potenciales electromagn\'eticos}

Como $\nabla\cdot\vb B=0$, el campo magn\'etico puede escribirse como el
rotacional de un potencial vectorial:
\begin{equation}
  \vb B=\nabla\times\vb A.
\end{equation}
Esta elecci\'on satisface autom\'aticamente la ley de Gauss magn\'etica, pues
\begin{equation}
  \nabla\cdot\vb B
  =
  \nabla\cdot(\nabla\times\vb A)
  =
  0.
\end{equation}
Sustituyendo $\vb B=\nabla\times\vb A$ en la ley de Faraday,
\begin{equation}
  \nabla\times\vb E
  =
  -\frac{\partial}{\partial t}
  (\nabla\times\vb A),
\end{equation}
queda
\begin{equation}
  \nabla\times
  \left(
  \vb E+\frac{\partial\vb A}{\partial t}
  \right)
  =
  0.
\end{equation}
Un campo con rotacional nulo puede escribirse localmente como el gradiente de
un potencial escalar. Por convenci\'on se escribe
\begin{equation}
  \vb E+\frac{\partial\vb A}{\partial t}
  =
  -\nabla\Phi,
\end{equation}
de modo que
\begin{equation}
  \boxed{
  \vb E=-\nabla\Phi-\frac{\partial\vb A}{\partial t},
  \qquad
  \vb B=\nabla\times\vb A.
  }
\end{equation}
Con este ansatz las ecuaciones homog\'eneas de Maxwell quedan satisfechas por
construcci\'on. Las ecuaciones inhomog\'eneas, que contienen las fuentes
$\rho$ y $\vb J$, determinan las ecuaciones para $\Phi$ y $\vb A$.

Sustituyendo en la ley de Ampere-Maxwell,
\begin{equation}
  \nabla\times\vb B
  =
  \mu_0\vb J
  +
  \frac{1}{c^2}\frac{\partial\vb E}{\partial t},
\end{equation}
con $c^{-2}=\mu_0\varepsilon_0$, se obtiene
\begin{align}
  \nabla\times(\nabla\times\vb A)
  &=
  \mu_0\vb J
  +
  \frac{1}{c^2}
  \frac{\partial}{\partial t}
  \left(
  -\nabla\Phi-\frac{\partial\vb A}{\partial t}
  \right)\\
  &=
  \mu_0\vb J
  -
  \frac{1}{c^2}
  \nabla\frac{\partial\Phi}{\partial t}
  -
  \frac{1}{c^2}
  \frac{\partial^2\vb A}{\partial t^2}.
\end{align}
Usando la identidad
\begin{equation}
  \nabla\times(\nabla\times\vb A)
  =
  \nabla(\nabla\cdot\vb A)-\nabla^2\vb A,
\end{equation}
se llega a
\begin{equation}
  \nabla^2\vb A
  -
  \frac{1}{c^2}
  \frac{\partial^2\vb A}{\partial t^2}
  =
  -\mu_0\vb J
  +
  \nabla
  \left(
  \nabla\cdot\vb A
  +
  \frac{1}{c^2}\frac{\partial\Phi}{\partial t}
  \right).
\end{equation}
Por otro lado, usando la ley de Gauss,
\begin{align}
  \nabla\cdot\vb E
  &=
  \nabla\cdot
  \left(
  -\nabla\Phi-\frac{\partial\vb A}{\partial t}
  \right)\\
  &=
  -\nabla^2\Phi
  -
  \frac{\partial}{\partial t}
  (\nabla\cdot\vb A)
  =
  \frac{\rho}{\varepsilon_0}.
\end{align}
Por tanto,
\begin{equation}
  \nabla^2\Phi
  -
  \frac{1}{c^2}
  \frac{\partial^2\Phi}{\partial t^2}
  =
  -\frac{\rho}{\varepsilon_0}
  -
  \frac{\partial}{\partial t}
  \left(
  \nabla\cdot\vb A
  +
  \frac{1}{c^2}
  \frac{\partial\Phi}{\partial t}
  \right).
\end{equation}
As\'i, en general las ecuaciones para $\Phi$ y $\vb A$ est\'an acopladas.

Los potenciales no son \'unicos. Si se transforma
\begin{equation}
  \vb A'=\vb A+\nabla\Lambda,
  \qquad
  \Phi'=\Phi-\frac{\partial\Lambda}{\partial t},
\end{equation}
entonces los campos f\'isicos no cambian. En efecto,
\begin{equation}
  \vb B'
  =
  \nabla\times\vb A'
  =
  \nabla\times\vb A+\nabla\times(\nabla\Lambda)
  =
  \vb B,
\end{equation}
y
\begin{align}
  \vb E'
  &=
  -\nabla\Phi'
  -
  \frac{\partial\vb A'}{\partial t}\\
  &=
  -\nabla\Phi
  +
  \nabla\frac{\partial\Lambda}{\partial t}
  -
  \frac{\partial\vb A}{\partial t}
  -
  \frac{\partial}{\partial t}(\nabla\Lambda)
  =
  \vb E.
\end{align}
Esta libertad de gauge permite imponer condiciones convenientes sobre los
potenciales.

Si se escoge $\Lambda$ de forma que los nuevos potenciales satisfagan
\begin{equation}
  \nabla\cdot\vb A'
  +
  \frac{1}{c^2}\frac{\partial\Phi'}{\partial t}
  =
  0,
\end{equation}
entonces se obtiene el gauge de Lorenz. Para verlo, se sustituye la
transformaci\'on de gauge:
\begin{align}
  0
  &=
  \nabla\cdot(\vb A+\nabla\Lambda)
  +
  \frac{1}{c^2}
  \frac{\partial}{\partial t}
  \left(
  \Phi-\frac{\partial\Lambda}{\partial t}
  \right)\\
  &=
  \nabla\cdot\vb A
  +
  \nabla^2\Lambda
  +
  \frac{1}{c^2}\frac{\partial\Phi}{\partial t}
  -
  \frac{1}{c^2}
  \frac{\partial^2\Lambda}{\partial t^2}.
\end{align}
Por tanto, basta resolver
\begin{equation}
  \left(
  \nabla^2-\frac{1}{c^2}\frac{\partial^2}{\partial t^2}
  \right)\Lambda
  =
  -
  \left(
  \nabla\cdot\vb A
  +
  \frac{1}{c^2}\frac{\partial\Phi}{\partial t}
  \right).
\end{equation}
Bajo el gauge de Lorenz, las ecuaciones de los potenciales se desacoplan:
\begin{equation}
  \nabla^2\Phi
  -
  \frac{1}{c^2}
  \frac{\partial^2\Phi}{\partial t^2}
  =
  -\frac{\rho}{\varepsilon_0},
\end{equation}
\begin{equation}
  \nabla^2\vb A
  -
  \frac{1}{c^2}
  \frac{\partial^2\vb A}{\partial t^2}
  =
  -\mu_0\vb J.
\end{equation}
Este gauge es muy \'util en problemas relativistas porque la condici\'on es
invariante de Lorentz.

En \'optica cu\'antica de cavidades se usa con frecuencia el gauge de
Coulomb,
\begin{equation}
  \nabla\cdot\vb A=0.
\end{equation}
Con esta elecci\'on,
\begin{equation}
  \nabla^2\Phi=-\frac{\rho}{\varepsilon_0}.
\end{equation}
En el vac\'io, donde $\rho=0$ y $\vb J=0$, se puede tomar $\Phi=0$, y el
potencial vectorial satisface la ecuaci\'on de onda
\begin{equation}
  \boxed{
  \nabla^2\vb A
  -
  \frac{1}{c^2}
  \frac{\partial^2\vb A}{\partial t^2}
  =
  0.
  }
\end{equation}
Resolver esta ecuaci\'on permite determinar $\vb E$ y $\vb B$ en cualquier
punto y en cualquier tiempo.

\section{Separaci\'on de variables en una cavidad}

Se considera un resonador de paredes conductoras. En el vac\'io y en gauge de
Coulomb se busca una soluci\'on separable del potencial vectorial:
\begin{equation}
  \vb A(\vb r,t)=\gamma\,q(t)\,\vb v(\vb r).
\end{equation}
Aqu\'i $q(t)$ contiene toda la dependencia temporal, $\vb v(\vb r)$ contiene
la forma espacial del modo, y $\gamma$ es una constante que se fija despu\'es
para que $q$ tenga unidades y normalizaci\'on convenientes.

Sustituyendo en la ecuaci\'on de onda,
\begin{equation}
  \nabla^2(\gamma q(t)\vb v(\vb r))
  -
  \frac{1}{c^2}
  \frac{\partial^2}{\partial t^2}
  (\gamma q(t)\vb v(\vb r))
  =
  0,
\end{equation}
se obtiene
\begin{equation}
  q(t)\nabla^2\vb v(\vb r)
  -
  \frac{1}{c^2}
  \ddot q(t)\vb v(\vb r)
  =
  0.
\end{equation}
Para cada componente $j=x,y,z$,
\begin{equation}
  \frac{\nabla^2 v_j(\vb r)}{v_j(\vb r)}
  =
  \frac{1}{c^2}
  \frac{\ddot q(t)}{q(t)}.
\end{equation}
El lado izquierdo s\'olo depende de la posici\'on, mientras que el derecho
s\'olo depende del tiempo. Para que la igualdad sea cierta para todo
$\vb r$ y todo $t$, ambos lados deben ser una constante. Se escoge
\begin{equation}
  \frac{\nabla^2 v_j(\vb r)}{v_j(\vb r)}
  =
  -k^2,
  \qquad
  \frac{1}{c^2}\frac{\ddot q(t)}{q(t)}
  =
  -k^2.
\end{equation}
As\'i se obtienen dos ecuaciones:
\begin{equation}
  \boxed{
  \nabla^2\vb v(\vb r)+k^2\vb v(\vb r)=0,
  }
\end{equation}
\begin{equation}
  \boxed{
  \ddot q(t)+\omega^2q(t)=0,
  \qquad
  \omega=ck.
  }
\end{equation}
La primera es la ecuaci\'on de Helmholtz para la estructura espacial del modo.
La segunda es la ecuaci\'on de un oscilador arm\'onico para la amplitud
temporal del modo. Esta es la primera se\~nal de que cada modo de cavidad se
comporta como un oscilador.

Las paredes conductoras imponen condiciones de frontera. En un conductor
perfecto, el campo el\'ectrico tangencial se anula en la superficie y el campo
magn\'etico normal tambi\'en:
\begin{equation}
  \vb E_{\parallel}=0,
  \qquad
  \vb B_{\perp}=0.
\end{equation}
Como en el vac\'io se toma $\Phi=0$,
\begin{equation}
  \vb E=-\frac{\partial\vb A}{\partial t}
  =
  -\gamma\dot q(t)\vb v(\vb r).
\end{equation}
Si $\uv{\parallel}(\vb r)$ es un vector tangente unitario a la frontera,
entonces
\begin{equation}
  E_\parallel
  =
  \uv{\parallel}\cdot\vb E\big|_F
  =
  -\gamma\dot q(t)
  \uv{\parallel}\cdot\vb v(\vb r)\big|_F
  =
  0.
\end{equation}
Para que esto se cumpla para todo tiempo,
\begin{equation}
  \boxed{
  \uv{\parallel}\cdot\vb v(\vb r)\big|_F=0.
  }
\end{equation}
Es decir, $\vb v$ debe ser normal a la superficie conductora.

La condici\'on magn\'etica se escribe con $\vb B=\nabla\times\vb A$:
\begin{equation}
  B_\perp
  =
  \uv{\perp}\cdot\vb B\big|_F
  =
  \gamma q(t)
  \uv{\perp}\cdot
  [\nabla\times\vb v(\vb r)]\big|_F
  =
  0.
\end{equation}
Por tanto,
\begin{equation}
  \boxed{
  \uv{\perp}\cdot
  [\nabla\times\vb v(\vb r)]\big|_F=0.
  }
\end{equation}
Estas dos condiciones no son independientes en el resonador rectangular.
Si $\vb v$ es normal a la superficie, la circulaci\'on de $\vb v$ sobre un lazo
infinitesimal tangente a la superficie es cero, porque $\dd\vb \ell$ es
tangente y $\vb v\cdot\dd\vb \ell=0$. Por el teorema de Stokes,
\begin{equation}
  \oint_C\vb v\cdot\dd\vb \ell
  =
  \int_S(\nabla\times\vb v)\cdot\uv{\perp}\,\dd S.
\end{equation}
Como el lado izquierdo se anula para un lazo tangente, la componente normal
del rotacional tambi\'en se anula en la frontera. Por eso la condici\'on
$B_\perp=0$ queda satisfecha una vez que se impone adecuadamente
$E_\parallel=0$.

En el resto del espacio se conserva la condici\'on de Coulomb. Con el ansatz,
\begin{equation}
  \nabla\cdot\vb A
  =
  \nabla\cdot(\gamma q(t)\vb v(\vb r))
  =
  \gamma q(t)\nabla\cdot\vb v(\vb r)
  =
  0.
\end{equation}
Por tanto,
\begin{equation}
  \boxed{\nabla\cdot\vb v(\vb r)=0.}
\end{equation}

\section{Modos de un resonador rectangular}

Se toma una cavidad rectangular de lados $L_x,L_y,L_z$. Sus caras se ubican
en $x=0$, $x=L_x$, $y=0$, $y=L_y$, $z=0$ y $z=L_z$. Se escribe
\begin{equation}
  \vb k=k_x\uv{x}+k_y\uv{y}+k_z\uv{z},
  \qquad
  \vb r=x\uv{x}+y\uv{y}+z\uv{z}.
\end{equation}
La condici\'on $\uv{\parallel}\cdot\vb v|_F=0$ implica, por ejemplo, que en
la cara $x=0$ las componentes tangenciales son $v_y$ y $v_z$, de modo que
\begin{equation}
  v_y(0,y,z)=0,
  \qquad
  v_z(0,y,z)=0.
\end{equation}
An\'alogamente, en $y=0$,
\begin{equation}
  v_x(x,0,z)=0,
  \qquad
  v_z(x,0,z)=0,
\end{equation}
y en $z=0$,
\begin{equation}
  v_x(x,y,0)=0,
  \qquad
  v_y(x,y,0)=0.
\end{equation}
Un ansatz que satisface estas condiciones es
\begin{align}
  v_x(x,y,z)
  &=
  N e_x\cos(k_xx)\sin(k_yy)\sin(k_zz),\\
  v_y(x,y,z)
  &=
  N e_y\sin(k_xx)\cos(k_yy)\sin(k_zz),\\
  v_z(x,y,z)
  &=
  N e_z\sin(k_xx)\sin(k_yy)\cos(k_zz).
\end{align}
El vector $\vb e=(e_x,e_y,e_z)$ fija la polarizaci\'on y $N$ fija la
normalizaci\'on.

La ecuaci\'on de Helmholtz se cumple componente por componente. Por ejemplo,
para $v_x$,
\begin{equation}
  \frac{\partial^2 v_x}{\partial x^2}
  =
  -k_x^2 v_x,
  \qquad
  \frac{\partial^2 v_x}{\partial y^2}
  =
  -k_y^2 v_x,
  \qquad
  \frac{\partial^2 v_x}{\partial z^2}
  =
  -k_z^2 v_x.
\end{equation}
Luego
\begin{equation}
  \nabla^2 v_x
  =
  -(k_x^2+k_y^2+k_z^2)v_x.
\end{equation}
Por tanto, la ecuaci\'on $\nabla^2\vb v+k^2\vb v=0$ exige
\begin{equation}
  k^2=k_x^2+k_y^2+k_z^2.
\end{equation}

El gauge de Coulomb impone
\begin{align}
  \nabla\cdot\vb v
  &=
  \partial_xv_x+\partial_yv_y+\partial_zv_z\\
  &=
  -N(e_xk_x+e_yk_y+e_zk_z)
  \sin(k_xx)\sin(k_yy)\sin(k_zz)
  =
  0.
\end{align}
Como la condici\'on debe cumplirse para todo $\vb r$,
\begin{equation}
  \boxed{
  e_xk_x+e_yk_y+e_zk_z
  =
  \vb e\cdot\vb k
  =
  0.
  }
\end{equation}
Esta es la condici\'on de transversalidad: la polarizaci\'on es ortogonal a la
direcci\'on de propagaci\'on. Para un $\vb k$ dado hay, en general, dos
direcciones de polarizaci\'on lineal independientes. Si alguna componente de
$\vb k$ se anula, el n\'umero de polarizaciones independientes puede reducirse
a una.

Las caras opuestas de la cavidad discretizan las componentes del vector de
onda. Las condiciones en $x=L_x$, $y=L_y$ y $z=L_z$ dan
\begin{equation}
  k_xL_x=\ell_x\pi,
  \qquad
  k_yL_y=\ell_y\pi,
  \qquad
  k_zL_z=\ell_z\pi,
\end{equation}
donde $\ell_x,\ell_y,\ell_z$ son enteros. Por tanto,
\begin{equation}
  \boxed{
  k_x=\ell_x\frac{\pi}{L_x},
  \qquad
  k_y=\ell_y\frac{\pi}{L_y},
  \qquad
  k_z=\ell_z\frac{\pi}{L_z}.
  }
\end{equation}
Esta discretizaci\'on de $\vb k$ viene de las condiciones de frontera de la
caja. No es todav\'ia la cuantizaci\'on del campo; la cuantizaci\'on del campo
ocurrir\'a al promover la amplitud $q(t)$ a operador.

Conviene verificar expl\'icitamente que las condiciones de $B_\perp$ se
satisfacen. En las caras $x=0$ y $x=L_x$, el vector normal es
$\pm\uv{x}$, de modo que se requiere
\begin{equation}
  (\nabla\times\vb v)_x
  =
  \frac{\partial v_z}{\partial y}
  -
  \frac{\partial v_y}{\partial z}
  =
  0
  \quad\text{en }x=0,L_x.
\end{equation}
Pero en esas caras se anulan $v_y$ y $v_z$ por la condici\'on tangencial
el\'ectrica, y el resultado se satisface autom\'aticamente. De forma similar,
\begin{equation}
  (\nabla\times\vb v)_y
  =
  \frac{\partial v_x}{\partial z}
  -
  \frac{\partial v_z}{\partial x}
  =
  0
  \quad\text{en }y=0,L_y,
\end{equation}
y
\begin{equation}
  (\nabla\times\vb v)_z
  =
  \frac{\partial v_y}{\partial x}
  -
  \frac{\partial v_x}{\partial y}
  =
  0
  \quad\text{en }z=0,L_z.
\end{equation}

\section{Ortogonalidad y normalizaci\'on}

El \'indice $\lambda$ representa cuatro datos: los enteros
$(\ell_x,\ell_y,\ell_z)$ y la polarizaci\'on. Se quiere construir una familia
de funciones modales que satisfaga
\begin{equation}
  \int_V \vb v_\lambda(\vb r)\cdot
  \vb v_{\lambda'}(\vb r)\,\dd^3r
  =
  \delta_{\lambda,\lambda'}.
\end{equation}
Si las polarizaciones son diferentes, la ortogonalidad se obtiene de
$\vb e_1\cdot\vb e_2=0$. Si la polarizaci\'on es la misma pero alg\'un entero
de modo es distinto, la ortogonalidad viene de las integrales de senos y
cosenos:
\begin{equation}
  \int_0^{L_i}
  \cos\left(\pi\ell_i\frac{x_i}{L_i}\right)
  \cos\left(\pi\ell_i'\frac{x_i}{L_i}\right)
  \dd x_i
  =
  \delta_{\ell_i,\ell_i'}\frac{L_i}{2},
\end{equation}
\begin{equation}
  \int_0^{L_i}
  \sin\left(\pi\ell_i\frac{x_i}{L_i}\right)
  \sin\left(\pi\ell_i'\frac{x_i}{L_i}\right)
  \dd x_i
  =
  \delta_{\ell_i,\ell_i'}\frac{L_i}{2}.
\end{equation}
Para la componente $x$, por ejemplo,
\begin{align}
  &\int_0^{L_z}\!\!\int_0^{L_y}\!\!\int_0^{L_x}
  v_{x,\lambda}v_{x,\lambda'}\,\dd x\,\dd y\,\dd z\\
  &=
  N_\lambda N_{\lambda'}
  e_{x,\lambda}e_{x,\lambda'}
  \int_0^{L_x}\cos(k_xx)\cos(k_x'x)\dd x
  \int_0^{L_y}\sin(k_yy)\sin(k_y'y)\dd y
  \int_0^{L_z}\sin(k_zz)\sin(k_z'z)\dd z.
\end{align}
Usando las relaciones anteriores queda
\begin{equation}
  N_\lambda N_{\lambda'}
  e_{x,\lambda}e_{x,\lambda'}
  \frac{L_xL_yL_z}{8}
  \delta_{\ell_x,\ell_x'}
  \delta_{\ell_y,\ell_y'}
  \delta_{\ell_z,\ell_z'}.
\end{equation}
Sumando las tres componentes y usando
\begin{equation}
  e_x^2+e_y^2+e_z^2=1,
\end{equation}
se obtiene, para $\lambda=\lambda'$,
\begin{equation}
  \int_V \vb v_\lambda^2(\vb r)\,\dd^3r
  =
  N_\lambda^2\frac{V}{8},
  \qquad
  V=L_xL_yL_z.
\end{equation}
La normalizaci\'on unitaria se consigue con
\begin{equation}
  N_\lambda=\sqrt{\frac{8}{V}}.
\end{equation}

En una cavidad de forma arbitraria se puede resumir la normalizaci\'on
introduciendo un volumen efectivo $V_\lambda$ y funciones adimensionales
$\vb u_\lambda(\vb r)$:
\begin{equation}
  \vb v_\lambda(\vb r)
  =
  \frac{1}{\sqrt{V_\lambda}}\vb u_\lambda(\vb r).
\end{equation}
Entonces la ortogonalidad se escribe como
\begin{equation}
  \frac{1}{\sqrt{V_\lambda V_{\lambda'}}}
  \int
  \vb u_\lambda(\vb r)\cdot
  \vb u_{\lambda'}(\vb r)\,\dd^3r
  =
  \delta_{\lambda,\lambda'}.
\end{equation}
Esta forma ser\'a la que se use para escribir la expansi\'on del campo.

\section{Expansi\'on modal del campo}

Con la normalizaci\'on anterior se expande el potencial vectorial como
\begin{equation}
  \vb A(\vb r,t)
  =
  \sum_\lambda
  \frac{1}{\sqrt{\varepsilon_0V_\lambda}}
  q_\lambda(t)\,
  \vb u_\lambda(\vb r).
\end{equation}
La elecci\'on del factor $1/\sqrt{\varepsilon_0V_\lambda}$ no es arbitraria:
se escoge para que la energ\'ia del campo tome exactamente la forma de una
suma de osciladores arm\'onicos con coordenada $q_\lambda$.

Como en el vac\'io $\Phi=0$,
\begin{equation}
  \vb E(\vb r,t)
  =
  -\frac{\partial\vb A}{\partial t}
  =
  -\sum_\lambda
  \frac{1}{\sqrt{\varepsilon_0V_\lambda}}
  \dot q_\lambda(t)\,
  \vb u_\lambda(\vb r).
\end{equation}
Adem\'as,
\begin{equation}
  \vb H(\vb r,t)
  =
  \frac{1}{\mu_0}\vb B(\vb r,t)
  =
  \frac{1}{\mu_0}\nabla\times\vb A(\vb r,t),
\end{equation}
de modo que
\begin{equation}
  \vb H(\vb r,t)
  =
  \sum_\lambda
  \frac{1}{\mu_0\sqrt{\varepsilon_0V_\lambda}}
  q_\lambda(t)
  \nabla\times\vb u_\lambda(\vb r).
\end{equation}
As\'i, el campo el\'ectrico es proporcional a $\dot q_\lambda$, mientras que
el campo magn\'etico es proporcional a $q_\lambda$. Esta correspondencia ser\'a
la raz\'on por la que, al cuantizar, $\vb E$ y $\vb B$ quedan asociados a
variables conjugadas.

\section{Energ\'ia del campo y teorema de Poynting}

Antes de sustituir los modos en la energ\'ia, conviene recordar de d\'onde
viene la expresi\'on de la energ\'ia electromagn\'etica. En presencia de cargas
y corrientes, la fuerza de Lorentz sobre una carga $q$ que se mueve con
velocidad $\vb v$ es
\begin{equation}
  \vb F=q(\vb E+\vb v\times\vb B).
\end{equation}
La potencia entregada por el campo a la carga es
\begin{equation}
  \frac{\dd W}{\dd t}
  =
  \vb F\cdot\vb v
  =
  q\vb E\cdot\vb v,
\end{equation}
porque $(\vb v\times\vb B)\cdot\vb v=0$. Si se tiene una densidad de carga
$\rho(\vb r,t)$ y una densidad de corriente
\begin{equation}
  \vb J(\vb r,t)=\rho(\vb r,t)\vb v(\vb r,t),
\end{equation}
entonces en un elemento de volumen $\dd^3r$ la carga es
\begin{equation}
  \dd q=\rho(\vb r,t)\dd^3r.
\end{equation}
La potencia total entregada a las cargas en un volumen $V$ es
\begin{equation}
  \frac{\dd W}{\dd t}
  =
  \int_V \vb E(\vb r,t)\cdot\vb J(\vb r,t)\,\dd^3r.
\end{equation}

Usando la ley de Ampere-Maxwell escrita como
\begin{equation}
  \nabla\times\vb H
  =
  \vb J+\frac{\partial\vb D}{\partial t},
\end{equation}
y en el vac\'io $\vb D=\varepsilon_0\vb E$, se obtiene
\begin{equation}
  \vb J
  =
  \nabla\times\vb H
  -
  \varepsilon_0\frac{\partial\vb E}{\partial t}.
\end{equation}
Entonces
\begin{align}
  \frac{\dd W}{\dd t}
  &=
  \int_V
  \vb E\cdot
  \left(
  \nabla\times\vb H
  -
  \varepsilon_0\frac{\partial\vb E}{\partial t}
  \right)
  \dd^3r\\
  &=
  \int_V
  \left[
  (\nabla\times\vb H)\cdot\vb E
  -
  \varepsilon_0
  \frac{\partial\vb E}{\partial t}\cdot\vb E
  \right]\dd^3r.
\end{align}
Se usa la identidad vectorial
\begin{equation}
  \nabla\cdot(\vb E\times\vb H)
  =
  \vb H\cdot(\nabla\times\vb E)
  -
  \vb E\cdot(\nabla\times\vb H).
\end{equation}
Por tanto,
\begin{equation}
  (\nabla\times\vb H)\cdot\vb E
  =
  \vb H\cdot(\nabla\times\vb E)
  -
  \nabla\cdot(\vb E\times\vb H).
\end{equation}
Con la ley de Faraday,
\begin{equation}
  \nabla\times\vb E
  =
  -\frac{\partial\vb B}{\partial t}
  =
  -\mu_0\frac{\partial\vb H}{\partial t},
\end{equation}
se obtiene
\begin{equation}
  \vb H\cdot(\nabla\times\vb E)
  =
  -\mu_0\vb H\cdot\frac{\partial\vb H}{\partial t}
  =
  -\frac{\mu_0}{2}
  \frac{\partial}{\partial t}H^2.
\end{equation}
De forma an\'aloga,
\begin{equation}
  -\varepsilon_0
  \frac{\partial\vb E}{\partial t}\cdot\vb E
  =
  -\frac{\varepsilon_0}{2}
  \frac{\partial}{\partial t}E^2.
\end{equation}
Por tanto,
\begin{equation}
  \frac{\dd W}{\dd t}
  =
  -\int_V
  \frac{\partial}{\partial t}
  \left(
  \frac{\varepsilon_0}{2}E^2
  +
  \frac{\mu_0}{2}H^2
  \right)\dd^3r
  -
  \int_V\nabla\cdot(\vb E\times\vb H)\,\dd^3r.
\end{equation}
Aplicando el teorema de Gauss al \'ultimo t\'ermino,
\begin{equation}
  \int_V\nabla\cdot(\vb E\times\vb H)\,\dd^3r
  =
  \oint_S(\vb E\times\vb H)\cdot\dd\vb S.
\end{equation}
As\'i se obtiene la ecuaci\'on integral de Poynting:
\begin{equation}
  \int_V \vb E\cdot\vb J\,\dd^3r
  +
  \int_V
  \frac{\partial}{\partial t}
  \left(
  \frac{\varepsilon_0}{2}E^2
  +
  \frac{\mu_0}{2}H^2
  \right)\dd^3r
  =
  -\oint_S(\vb E\times\vb H)\cdot\dd\vb S.
\end{equation}
En forma local,
\begin{equation}
  \boxed{
  \frac{\partial u}{\partial t}
  +
  \nabla\cdot\vb S
  =
  -\vb E\cdot\vb J,
  }
\end{equation}
donde
\begin{equation}
  u=
  \frac{\varepsilon_0}{2}E^2
  +
  \frac{\mu_0}{2}H^2,
  \qquad
  \vb S=\vb E\times\vb H.
\end{equation}
El primer t\'ermino es la variaci\'on local de energ\'ia almacenada en el
campo, el segundo es el flujo de energ\'ia, y el lado derecho es la potencia
transferida a las cargas. En una cavidad conductora sin cargas ni corrientes,
no hay flujo neto de energ\'ia hacia afuera y la energ\'ia total del campo se
conserva:
\begin{equation}
  \mathcal H
  =
  \int_V
  \left(
  \frac{\varepsilon_0}{2}E^2
  +
  \frac{\mu_0}{2}H^2
  \right)\dd^3r.
\end{equation}

\section{Reducci\'on a osciladores arm\'onicos}

Sustituyendo la expansi\'on modal en la energ\'ia,
\begin{align}
  \mathcal H
  &=
  \frac{\varepsilon_0}{2}
  \int
  \vb E^2\,\dd^3r
  +
  \frac{\mu_0}{2}
  \int
  \vb H^2\,\dd^3r\\
  &=
  \frac{1}{2}
  \sum_{\lambda,\lambda'}
  \frac{\dot q_\lambda\dot q_{\lambda'}}
  {\sqrt{V_\lambda V_{\lambda'}}}
  \int
  \vb u_\lambda\cdot\vb u_{\lambda'}\,\dd^3r\\
  &\quad+
  \frac{1}{2\mu_0\varepsilon_0}
  \sum_{\lambda,\lambda'}
  \frac{q_\lambda q_{\lambda'}}
  {\sqrt{V_\lambda V_{\lambda'}}}
  \int
  [\nabla\times\vb u_\lambda]\cdot
  [\nabla\times\vb u_{\lambda'}]\,\dd^3r.
\end{align}
El primer t\'ermino se simplifica por ortogonalidad:
\begin{equation}
  \frac{1}{2}
  \sum_{\lambda,\lambda'}
  \dot q_\lambda\dot q_{\lambda'}
  \delta_{\lambda,\lambda'}
  =
  \frac{1}{2}\sum_\lambda \dot q_\lambda^2.
\end{equation}
Para el t\'ermino magn\'etico se usa
\begin{equation}
  \nabla\cdot[\vb u_\lambda\times(\nabla\times\vb u_{\lambda'})]
  =
  [\nabla\times\vb u_\lambda]\cdot
  [\nabla\times\vb u_{\lambda'}]
  -
  \vb u_\lambda\cdot
  [\nabla\times(\nabla\times\vb u_{\lambda'})].
\end{equation}
Al integrar en el volumen, el primer t\'ermino del lado derecho produce una
integral de superficie. En la frontera conductora, $\vb E$ y $\vb B$ son
ortogonales de manera tal que el vector asociado queda tangente a la
superficie, mientras que $\dd\vb S$ es normal. Por eso la contribuci\'on de
superficie se anula. Queda
\begin{equation}
  \int
  [\nabla\times\vb u_\lambda]\cdot
  [\nabla\times\vb u_{\lambda'}]\,\dd^3r
  =
  \int
  \vb u_\lambda\cdot
  [\nabla\times(\nabla\times\vb u_{\lambda'})]\,\dd^3r.
\end{equation}
Usando
\begin{equation}
  \nabla\times(\nabla\times\vb u_{\lambda'})
  =
  \nabla(\nabla\cdot\vb u_{\lambda'})
  -
  \nabla^2\vb u_{\lambda'},
\end{equation}
la condici\'on de Coulomb $\nabla\cdot\vb u_{\lambda'}=0$ y la ecuaci\'on de
Helmholtz $\nabla^2\vb u_{\lambda'}+k_{\lambda'}^2\vb u_{\lambda'}=0$,
se obtiene
\begin{equation}
  \nabla\times(\nabla\times\vb u_{\lambda'})
  =
  k_{\lambda'}^2\vb u_{\lambda'}.
\end{equation}
Por tanto,
\begin{equation}
  \frac{1}{\sqrt{V_\lambda V_{\lambda'}}}
  \int
  \vb u_\lambda\cdot
  [\nabla\times(\nabla\times\vb u_{\lambda'})]\,\dd^3r
  =
  k_{\lambda'}^2\delta_{\lambda,\lambda'}.
\end{equation}
Como $c^2=1/(\mu_0\varepsilon_0)$ y $\omega_\lambda=ck_\lambda$,
la energ\'ia queda
\begin{equation}
  \boxed{
  \mathcal H
  =
  \sum_\lambda
  \left(
  \frac{1}{2}\dot q_\lambda^2
  +
  \frac{1}{2}\omega_\lambda^2q_\lambda^2
  \right).
  }
\end{equation}
La energ\'ia del campo electromagn\'etico en un resonador es la suma de las
energ\'ias de osciladores arm\'onicos independientes, uno por cada modo
$\lambda$.

La coordenada $q_\lambda$ se interpreta entonces como la coordenada normal del
modo y
\begin{equation}
  p_\lambda=\dot q_\lambda
\end{equation}
como su momento conjugado. Las ecuaciones de Hamilton son
\begin{equation}
  \dot q_\lambda
  =
  \frac{\partial\mathcal H_\lambda}{\partial p_\lambda}
  =
  p_\lambda,
  \qquad
  \dot p_\lambda
  =
  -\frac{\partial\mathcal H_\lambda}{\partial q_\lambda}
  =
  -\omega_\lambda^2q_\lambda,
\end{equation}
y conducen a
\begin{equation}
  \ddot q_\lambda+\omega_\lambda^2q_\lambda=0.
\end{equation}
Esto coincide con la ecuaci\'on temporal obtenida por separaci\'on de
variables.

\section{Cuantizaci\'on de las amplitudes normales}

La cuantizaci\'on consiste en promover $q_\lambda$ y $p_\lambda$ a operadores:
\begin{equation}
  q_\lambda\longrightarrow\hat q_\lambda,
  \qquad
  p_\lambda\longrightarrow\hat p_\lambda,
\end{equation}
e imponer
\begin{equation}
  [\hat q_\lambda,\hat p_{\lambda'}]
  =
  \ii\hbar\,\delta_{\lambda,\lambda'},
  \qquad
  [\hat q_\lambda,\hat q_{\lambda'}]=
  [\hat p_\lambda,\hat p_{\lambda'}]=0.
\end{equation}
Para cada modo se introducen operadores escalares de aniquilaci\'on y
creaci\'on:
\begin{equation}
  \hat a_\lambda
  =
  \frac{1}{\sqrt{2\hbar\omega_\lambda}}
  (\omega_\lambda\hat q_\lambda+\ii\hat p_\lambda),
\end{equation}
\begin{equation}
  \hat a_\lambda^\dagger
  =
  \frac{1}{\sqrt{2\hbar\omega_\lambda}}
  (\omega_\lambda\hat q_\lambda-\ii\hat p_\lambda).
\end{equation}
Sumando y restando estas definiciones se obtienen las relaciones inversas:
\begin{equation}
  \hat q_\lambda
  =
  \sqrt{\frac{\hbar}{2\omega_\lambda}}
  \left(
  \hat a_\lambda+\hat a_\lambda^\dagger
  \right),
\end{equation}
\begin{equation}
  \hat p_\lambda
  =
  -\ii\sqrt{\frac{\hbar\omega_\lambda}{2}}
  \left(
  \hat a_\lambda-\hat a_\lambda^\dagger
  \right).
\end{equation}
El conmutador de $\hat a_\lambda$ y $\hat a_{\lambda'}^\dagger$ se obtiene
directamente:
\begin{align}
  [\hat a_\lambda,\hat a_{\lambda'}^\dagger]
  &=
  \frac{1}{\sqrt{2\hbar\omega_\lambda}}
  \frac{1}{\sqrt{2\hbar\omega_{\lambda'}}}
  [\omega_\lambda\hat q_\lambda+\ii\hat p_\lambda,
  \omega_{\lambda'}\hat q_{\lambda'}-\ii\hat p_{\lambda'}].
\end{align}
Los conmutadores $[\hat q,\hat q]$ y $[\hat p,\hat p]$ se anulan, por lo que
quedan s\'olo los t\'erminos cruzados. Para $\lambda=\lambda'$,
\begin{align}
  [\hat a_\lambda,\hat a_{\lambda}^\dagger]
  &=
  \frac{1}{2\hbar\omega_\lambda}
  \left(
  -\ii\omega_\lambda[\hat q_\lambda,\hat p_\lambda]
  +
  \ii\omega_\lambda[\hat p_\lambda,\hat q_\lambda]
  \right)\\
  &=
  \frac{1}{2\hbar\omega_\lambda}
  \left(
  -\ii\omega_\lambda(\ii\hbar)
  +
  \ii\omega_\lambda(-\ii\hbar)
  \right)
  =
  1.
\end{align}
En general,
\begin{equation}
  \boxed{
  [\hat a_\lambda,\hat a_{\lambda'}^\dagger]
  =
  \delta_{\lambda,\lambda'},
  \qquad
  [\hat a_\lambda,\hat a_{\lambda'}]=
  [\hat a_\lambda^\dagger,\hat a_{\lambda'}^\dagger]=0.
  }
\end{equation}

El hamiltoniano de un modo se escribe como
\begin{equation}
  \hat{\mathcal H}_\lambda
  =
  \frac{1}{2}\hat p_\lambda^2
  +
  \frac{1}{2}\omega_\lambda^2\hat q_\lambda^2.
\end{equation}
Sustituyendo las expresiones de $\hat q_\lambda$ y $\hat p_\lambda$,
\begin{align}
  \hat{\mathcal H}_\lambda
  &=
  \frac{1}{2}
  \left[
  -\ii\sqrt{\frac{\hbar\omega_\lambda}{2}}
  (\hat a_\lambda-\hat a_\lambda^\dagger)
  \right]^2
  +
  \frac{1}{2}\omega_\lambda^2
  \left[
  \sqrt{\frac{\hbar}{2\omega_\lambda}}
  (\hat a_\lambda+\hat a_\lambda^\dagger)
  \right]^2\\
  &=
  -\frac{\hbar\omega_\lambda}{4}
  (\hat a_\lambda-\hat a_\lambda^\dagger)^2
  +
  \frac{\hbar\omega_\lambda}{4}
  (\hat a_\lambda+\hat a_\lambda^\dagger)^2.
\end{align}
Expandiendo,
\begin{align}
  (\hat a+\hat a^\dagger)^2
  &=
  \hat a^2+\hat a\hat a^\dagger
  +
  \hat a^\dagger\hat a
  +
  \hat a^{\dagger 2},\\
  (\hat a-\hat a^\dagger)^2
  &=
  \hat a^2-\hat a\hat a^\dagger
  -
  \hat a^\dagger\hat a
  +
  \hat a^{\dagger 2}.
\end{align}
Los t\'erminos $\hat a^2$ y $\hat a^{\dagger2}$ se cancelan, y queda
\begin{equation}
  \hat{\mathcal H}_\lambda
  =
  \frac{\hbar\omega_\lambda}{2}
  \left(
  \hat a_\lambda\hat a_\lambda^\dagger
  +
  \hat a_\lambda^\dagger\hat a_\lambda
  \right).
\end{equation}
Usando
\begin{equation}
  \hat a_\lambda\hat a_\lambda^\dagger
  =
  \hat a_\lambda^\dagger\hat a_\lambda+1,
\end{equation}
se obtiene
\begin{equation}
  \hat{\mathcal H}_\lambda
  =
  \hbar\omega_\lambda
  \left(
  \hat a_\lambda^\dagger\hat a_\lambda+\frac{1}{2}
  \right).
\end{equation}
Por tanto, el hamiltoniano del campo en la cavidad es
\begin{equation}
  \boxed{
  \hat{\mathcal H}
  =
  \sum_\lambda
  \hbar\omega_\lambda
  \hat a_\lambda^\dagger\hat a_\lambda
  +
  \mathcal H_0,
  \qquad
  \mathcal H_0=
  \sum_\lambda\frac{1}{2}\hbar\omega_\lambda.
  }
\end{equation}
El t\'ermino $\mathcal H_0$ viene de la relaci\'on de conmutaci\'on. No
contiene operadores de n\'umero, pero suma la energ\'ia de punto cero de todos
los modos.

\section{Energ\'ia de punto cero y efecto Casimir}

La frecuencia del modo $\lambda$ en una cavidad rectangular es
\begin{equation}
  \omega_\lambda
  =
  c|\vb k_\lambda|
  =
  c
  \left[
  \frac{\ell_x^2\pi^2}{L_x^2}
  +
  \frac{\ell_y^2\pi^2}{L_y^2}
  +
  \frac{\ell_z^2\pi^2}{L_z^2}
  \right]^{1/2}.
\end{equation}
Aunque la suma de energ\'ias de punto cero sea infinita, sus diferencias
pueden depender de la geometr\'ia y producir efectos observables. Para ver
esto se consideran dos placas paralelas de \'area $L^2$ separadas una
distancia $a$. Se toma
\begin{equation}
  L_x=a,
  \qquad
  L_y=L_z=L.
\end{equation}
Entonces
\begin{equation}
  \omega_\lambda
  =
  c
  \left[
  \frac{\ell_x^2\pi^2}{a^2}
  +
  (\ell_y^2+\ell_z^2)\frac{\pi^2}{L^2}
  \right]^{1/2}.
\end{equation}
Por cada $\omega_\lambda$ hay dos polarizaciones, excepto cuando alguna
componente del vector de onda se anula. En particular, el modo con
$\ell_x=0$ contribuye con una sola direcci\'on de polarizaci\'on.

La energ\'ia de punto cero dentro de la caja se escribe como
\begin{equation}
  \mathcal H_0^{\mathrm{box}}
  =
  \frac{\hbar c}{2}
  \sum_{\mathrm{pol}}
  \sum_{\ell_x,\ell_y,\ell_z=0}^{\infty}
  \left[
  \frac{\ell_x^2\pi^2}{a^2}
  +
  (\ell_y^2+\ell_z^2)\frac{\pi^2}{L^2}
  \right]^{1/2}.
\end{equation}
Si $L\gg a$, los modos en las direcciones $y$ y $z$ son muy densos. Se
aproximan las sumas sobre $\ell_y$ y $\ell_z$ por integrales. Definiendo
\begin{equation}
  \xi=\frac{a}{L}\ell_y,
  \qquad
  \zeta=\frac{a}{L}\ell_z,
\end{equation}
se tiene $\dd \ell_y\dd\ell_z=(L^2/a^2)\dd\xi\dd\zeta$, y
\begin{equation}
  \mathcal H_0^{\mathrm{box}}
  =
  \frac{\hbar c\pi}{2a}
  \frac{L^2}{a^2}
  \sum_{\mathrm{pol}}
  \sum_{\ell_x=0}^{\infty}
  \int_0^\infty\int_0^\infty
  (\ell_x^2+\xi^2+\zeta^2)^{1/2}
  \dd\xi\,\dd\zeta.
\end{equation}
Pasando a coordenadas polares en el primer cuadrante,
\begin{equation}
  \xi=\sqrt{u}\cos\varphi,
  \qquad
  \zeta=\sqrt{u}\sin\varphi,
\end{equation}
el jacobiano es
\begin{equation}
  J
  =
  \left|
  \begin{matrix}
  \frac{1}{2\sqrt u}\cos\varphi & -\sqrt u\sin\varphi\\
  \frac{1}{2\sqrt u}\sin\varphi & \sqrt u\cos\varphi
  \end{matrix}
  \right|
  =
  \frac{1}{2}.
\end{equation}
Como $\varphi$ va de $0$ a $\pi/2$, se obtiene
\begin{equation}
  \mathcal H_0^{\mathrm{box}}
  =
  \frac{\hbar c\pi^2}{8a^3}L^2
  \sum_{\mathrm{pol}}
  \sum_{\ell_x=0}^{\infty}
  \int_0^\infty
  (\ell_x^2+u)^{1/2}\dd u.
\end{equation}
Separando el t\'ermino $\ell_x=0$, que tiene una sola polarizaci\'on,
\begin{equation}
  \mathcal H_0^{\mathrm{box}}
  =
  \frac{\hbar c\pi^2}{8a^3}L^2
  \left[
  \int_0^\infty \sqrt u\,\dd u
  +
  2\sum_{\ell_x=1}^{\infty}
  \int_0^\infty
  \sqrt{\ell_x^2+u}\,\dd u
  \right].
\end{equation}
El primer t\'ermino y la suma son divergentes.

Para construir la referencia sin placas se deja que la separaci\'on tambi\'en
se vuelva continua. La energ\'ia del vac\'io libre es
\begin{equation}
  \mathcal H_0^{\mathrm{vac}}
  =
  \frac{\hbar c\pi^2}{8a^3}L^2
  2\int_0^\infty\int_0^\infty
  (\ell_x^2+u)^{1/2}\dd \ell_x\,\dd u.
\end{equation}
La cantidad f\'isica es la diferencia entre la energ\'ia con placas y la
energ\'ia del vac\'io libre. Por unidad de \'area,
\begin{equation}
  v
  =
  \frac{1}{L^2}
  \left(
  \mathcal H_0^{\mathrm{box}}
  -
  \mathcal H_0^{\mathrm{vac}}
  \right)
  =
  \frac{\hbar c\pi^2}{4a^3}\,\tilde I,
\end{equation}
donde
\begin{equation}
  \tilde I
  =
  \frac{1}{2}I(0)
  +
  \sum_{\ell_x=1}^{\infty}I(\ell_x)
  -
  \int_0^\infty I(\ell_x)\dd\ell_x,
\end{equation}
con
\begin{equation}
  I(\lambda)
  =
  \int_0^\infty
  \sqrt{\lambda^2+u}\,\dd u.
\end{equation}
Para aplicar Euler-Maclaurin, la suma debe empezar en cero. Como
\begin{equation}
  \sum_{\ell_x=1}^{\infty}I(\ell_x)
  =
  \sum_{\ell_x=0}^{\infty}I(\ell_x)-I(0),
\end{equation}
queda
\begin{equation}
  \tilde I
  =
  -\frac{1}{2}I(0)
  +
  \sum_{\ell_x=0}^{\infty}I(\ell_x)
  -
  \int_0^\infty I(\ell_x)\dd\ell_x.
\end{equation}
La f\'ormula de Euler-Maclaurin dice que
\begin{equation}
  \sum_{n=a}^{b}F(n)
  =
  \int_a^bF(x)\dd x
  +
  \frac{1}{2}[F(a)+F(b)]
  +
  \sum_{k=1}^{p}
  \frac{B_{2k}}{(2k)!}
  [F^{(2k-1)}(b)-F^{(2k-1)}(a)]
  +
  R_p,
\end{equation}
donde $B_{2k}$ son los n\'umeros de Bernoulli. Aqu\'i
$B_2=1/6$ y $B_4=-1/30$.

La integral $I(\lambda)$ se eval\'ua formalmente haciendo
$v=\lambda^2+u$:
\begin{equation}
  I(\lambda)
  =
  \int_{\lambda^2}^{\infty}\sqrt v\,\dd v
  =
  \left.
  \frac{2}{3}v^{3/2}
  \right|_{\lambda^2}^{\infty}
  =
  \infty-\frac{2}{3}\lambda^3.
\end{equation}
La parte infinita no depende de $\lambda$ en las derivadas finitas, por lo que
\begin{equation}
  I'(\lambda)=-2\lambda^2,
  \qquad
  I''(\lambda)=-4\lambda,
  \qquad
  I'''(\lambda)=-4,
  \qquad
  I^{(4)}=I^{(5)}=\cdots=0.
\end{equation}
Los t\'erminos divergentes comunes se cancelan al restar el vac\'io libre, y
el primer t\'ermino finito que queda es
\begin{equation}
  \tilde I
  =
  \frac{B_4}{4!}
  [I'''(\infty)-I'''(0)]
  =
  -\frac{1}{30}\frac{1}{24}[0-(-4)]
  =
  -\frac{1}{180}.
\end{equation}
Entonces
\begin{equation}
  \boxed{
  v(a)
  =
  -\frac{\hbar c\pi^2}{720a^3}.
  }
\end{equation}
Para placas de \'area $A$, la energ\'ia potencial es
\begin{equation}
  \boxed{
  V(a)
  =
  -\frac{\hbar c\pi^2}{720}\frac{A}{a^3}.
  }
\end{equation}
La fuerza por unidad de \'area se obtiene derivando:
\begin{equation}
  \boxed{
  \frac{F}{A}
  =
  -\frac{\dd v}{\dd a}
  =
  -\frac{\hbar c\pi^2}{240a^4}.
  }
\end{equation}
El signo negativo indica que la fuerza apunta hacia separaciones menores:
las placas se atraen. Esta fuerza es de origen cu\'antico, pues proviene de la
dependencia geom\'etrica de la energ\'ia de punto cero.

\section{Operadores del potencial vectorial y de los campos}

Despu\'es de cuantizar, la expansi\'on del potencial vectorial se convierte en
\begin{equation}
  \hat{\vb A}(\vb r,t)
  =
  \sum_\lambda
  \frac{1}{\sqrt{\varepsilon_0V_\lambda}}
  \hat q_\lambda(t)\vb u_\lambda(\vb r).
\end{equation}
Usando
\begin{equation}
  \hat q_\lambda
  =
  \sqrt{\frac{\hbar}{2\omega_\lambda}}
  (\hat a_\lambda+\hat a_\lambda^\dagger),
\end{equation}
se obtiene
\begin{equation}
  \hat{\vb A}(\vb r,t)
  =
  \sum_\lambda
  \sqrt{
  \frac{\hbar}{2\varepsilon_0V_\lambda\omega_\lambda}
  }
  \vb u_\lambda(\vb r)
  [\hat a_\lambda(t)+\hat a_\lambda^\dagger(t)].
\end{equation}

La dependencia temporal en la imagen de Heisenberg se obtiene de
\begin{equation}
  \dot{\hat a}_\lambda
  =
  \frac{\ii}{\hbar}
  [\hat{\mathcal H},\hat a_\lambda].
\end{equation}
Usando
\begin{equation}
  \hat{\mathcal H}
  =
  \sum_{\lambda'}
  \hbar\omega_{\lambda'}
  \hat a_{\lambda'}^\dagger\hat a_{\lambda'}
  +
  \mathcal H_0,
\end{equation}
y que $\mathcal H_0$ conmuta con todo, se calcula
\begin{align}
  [\hat a_{\lambda'}^\dagger\hat a_{\lambda'},\hat a_\lambda]
  &=
  \hat a_{\lambda'}^\dagger
  [\hat a_{\lambda'},\hat a_\lambda]
  +
  [\hat a_{\lambda'}^\dagger,\hat a_\lambda]\hat a_{\lambda'}\\
  &=
  -\delta_{\lambda,\lambda'}\hat a_{\lambda'}.
\end{align}
Por tanto,
\begin{equation}
  \dot{\hat a}_\lambda
  =
  -\ii\omega_\lambda\hat a_\lambda,
  \qquad
  \hat a_\lambda(t)
  =
  \hat a_\lambda(0)\ee^{-\ii\omega_\lambda t}.
\end{equation}
Similarmente,
\begin{equation}
  \hat a_\lambda^\dagger(t)
  =
  \hat a_\lambda^\dagger(0)\ee^{\ii\omega_\lambda t}.
\end{equation}
Definiendo
\begin{equation}
  A_\lambda
  =
  \sqrt{\frac{\hbar}{\varepsilon_0V_\lambda\omega_\lambda}},
\end{equation}
el operador potencial vectorial queda
\begin{equation}
  \boxed{
  \hat{\vb A}(\vb r,t)
  =
  \sum_\lambda
  \frac{A_\lambda}{\sqrt2}\,
  \vb u_\lambda(\vb r)
  \left[
  \hat a_\lambda(0)\ee^{-\ii\omega_\lambda t}
  +
  \hat a_\lambda^\dagger(0)\ee^{\ii\omega_\lambda t}
  \right].
  }
\end{equation}
La parte de frecuencia positiva contiene operadores de aniquilaci\'on,
\begin{equation}
  \hat{\vb A}^{(+)}(\vb r,t)
  =
  \frac{1}{\sqrt2}
  \sum_\lambda
  A_\lambda
  \vb u_\lambda(\vb r)
  \hat a_\lambda(0)\ee^{-\ii\omega_\lambda t},
\end{equation}
y la parte de frecuencia negativa contiene operadores de creaci\'on,
\begin{equation}
  \hat{\vb A}^{(-)}(\vb r,t)
  =
  \frac{1}{\sqrt2}
  \sum_\lambda
  A_\lambda
  \vb u_\lambda(\vb r)
  \hat a_\lambda^\dagger(0)\ee^{\ii\omega_\lambda t}.
\end{equation}
Por construcci\'on,
\begin{equation}
  \hat{\vb A}=\hat{\vb A}^{(+)}+\hat{\vb A}^{(-)}.
\end{equation}

El campo el\'ectrico operatorial es
\begin{equation}
  \hat{\vb E}(\vb r,t)
  =
  -\frac{\partial\hat{\vb A}}{\partial t}.
\end{equation}
Derivando,
\begin{equation}
  \boxed{
  \hat{\vb E}(\vb r,t)
  =
  \sum_\lambda
  \frac{\ii E_\lambda}{\sqrt2}
  \vb u_\lambda(\vb r)
  \left[
  \hat a_\lambda(0)\ee^{-\ii\omega_\lambda t}
  -
  \hat a_\lambda^\dagger(0)\ee^{\ii\omega_\lambda t}
  \right],
  }
\end{equation}
donde
\begin{equation}
  E_\lambda=\omega_\lambda A_\lambda
  =
  \sqrt{\frac{\hbar\omega_\lambda}{\varepsilon_0V_\lambda}}.
\end{equation}
La descomposici\'on en frecuencias es
\begin{equation}
  \hat{\vb E}^{(+)}(\vb r,t)
  =
  \frac{\ii}{\sqrt2}
  \sum_\lambda
  E_\lambda
  \vb u_\lambda(\vb r)
  \hat a_\lambda(0)\ee^{-\ii\omega_\lambda t},
\end{equation}
\begin{equation}
  \hat{\vb E}^{(-)}(\vb r,t)
  =
  -\frac{\ii}{\sqrt2}
  \sum_\lambda
  E_\lambda
  \vb u_\lambda(\vb r)
  \hat a_\lambda^\dagger(0)\ee^{\ii\omega_\lambda t}.
\end{equation}

El campo magn\'etico se obtiene con
\begin{equation}
  \hat{\vb B}(\vb r,t)=\nabla\times\hat{\vb A}(\vb r,t),
\end{equation}
por lo que
\begin{equation}
  \boxed{
  \hat{\vb B}(\vb r,t)
  =
  \sum_\lambda
  \frac{A_\lambda}{\sqrt2}
  [\nabla\times\vb u_\lambda(\vb r)]
  \left[
  \hat a_\lambda(0)\ee^{-\ii\omega_\lambda t}
  +
  \hat a_\lambda^\dagger(0)\ee^{\ii\omega_\lambda t}
  \right].
  }
\end{equation}
Equivalentemente,
\begin{equation}
  \hat{\vb H}(\vb r,t)
  =
  \frac{1}{\mu_0}
  \hat{\vb B}(\vb r,t).
\end{equation}
Como $\hat{\vb E}$ contiene la variable conjugada de momento y
$\hat{\vb B}$ contiene la variable de posici\'on del oscilador, los campos
el\'ectrico y magn\'etico operatoriales no conmutan en general. Esta no
conmutatividad es la versi\'on de campo de la no conmutatividad entre
posici\'on y momento de cada oscilador arm\'onico.

\section{Estados de n\'umero del campo}

Si se omite la energ\'ia de punto cero mediante una renormalizaci\'on de la
energ\'ia, el hamiltoniano de un modo queda
\begin{equation}
  \hat{\mathcal H}_\lambda
  =
  \hbar\omega_\lambda\hat n_\lambda,
  \qquad
  \hat n_\lambda
  =
  \hat a_\lambda^\dagger\hat a_\lambda.
\end{equation}
Los eigenestados de $\hat n_\lambda$ satisfacen
\begin{equation}
  \hat n_\lambda\ket{n_\lambda}
  =
  n_\lambda\ket{n_\lambda},
\end{equation}
con
\begin{equation}
  n_\lambda=0,1,2,\ldots.
\end{equation}
Las excitaciones discretas $n_\lambda\hbar\omega_\lambda$ del modo son los
fotones. Para un solo modo, el estado $\ket{n_\lambda}$ contiene
$n_\lambda$ fotones en ese modo.

El campo completo contiene muchos modos independientes. Si cada modo est\'a
en un estado de n\'umero, el estado total se escribe como un producto tensorial
\begin{equation}
  \ket{\Psi}
  =
  \ket{n_1}\ket{n_2}\ket{n_3}\cdots\ket{n_\lambda}\cdots
  =
  \ket{n_1,n_2,n_3,\ldots,n_\lambda,\ldots}.
\end{equation}
Aqu\'i $n_\lambda$ indica el n\'umero de excitaciones en el modo $\lambda$.
El valor esperado de la energ\'ia, ignorando $\mathcal H_0$, es
\begin{align}
  \braket{\hat{\mathcal H}}
  &=
  \bra{n_1,n_2,\ldots}
  \sum_\lambda
  \hbar\omega_\lambda\hat n_\lambda
  \ket{n_1,n_2,\ldots}\\
  &=
  \sum_\lambda
  \hbar\omega_\lambda
  n_\lambda.
\end{align}
Se us\'o que los estados de n\'umero de modos distintos son ortogonales.

El estado de un modo no tiene que ser necesariamente un estado de n\'umero.
Puede ser una superposici\'on
\begin{equation}
  \ket{\psi_\lambda}
  =
  \sum_{n_\lambda}
  c_{n_\lambda}\ket{n_\lambda},
\end{equation}
y el estado de campo factorizado ser
\begin{equation}
  \ket{\Psi}
  =
  \ket{\psi_1}\ket{\psi_2}\ket{\psi_3}\cdots\ket{\psi_\lambda}\cdots.
\end{equation}
Algunos de estos estados son estados coherentes, estados apachurrados o
estados de cuadratura. Tambi\'en pueden existir estados no factorizables entre
modos, es decir, estados entrelazados del campo.

\end{document}
